\begin{textsample}{POS Dim 3 – persona\_gemini – Score 23.00 – t932\_gemini.txt}  \label{ex:f3_pos_041}
The World Health Organization classifies processed \textbf{meat} as carcinogenic and red \textbf{meat} as probably carcinogenic, raising concerns about current \textbf{meat} \textbf{consumption} \textbf{levels}.
The environmental impact of this \textbf{industry} is immense.
Animal \textbf{agriculture} is responsible for 80 % of Amazon deforestation and generates more \textbf{greenhouse} \textbf{gases} than the \textbf{transport} \textbf{sector}.
It also \textbf{uses} 80 % of the world 's soy for animal feed.
The \textbf{scale} of \textbf{production} involves breeding 1 \textbf{billion} pigs, 19 \textbf{billion} chickens, and 1.4 \textbf{billion} \textbf{cattle}.
This \textbf{level} of \textbf{production} has a significant resource \textbf{footprint}.
For instance, producing one person 's annual \textbf{supply} of \textbf{meat} and \textbf{dairy} requires 403,000 liters of water, an \textbf{amount} \textbf{equivalent} to 6,190 showers.
Additionally, the \textbf{system} relies heavily on medication ; in 2011, 8,481 \textbf{tons} of antibiotics were sold for \textbf{livestock} in the EU.
For significant \textbf{health}, ethical, and environmental benefits, we can reduce our weekly \textbf{meat} intake by \textbf{half} and choose to eat more local fruits and vegetables.

% matched lemmas: agriculture, amount, billion, cattle, consumption, dairy, equivalent, footprint, gas, greenhouse, half, health, industry, level, livestock, meat, production, scale, sector, supply, system, ton, transport, use
\end{textsample}
